\documentclass{article}
\usepackage[MeX]{polski}
\usepackage[utf8]{inputenc}
\usepackage{tikz}

\title{Moj artykul przy pomocy - \LaTeX}
\author{Adrian Kwasnik}


\begin{document}
\maketitle

\begin{abstract}
    Artykuł poświęcony jest fascynującemu zjawisku czarnych wielorybów, tajemniczym mieszkańcom głębin oceanicznych. Czarne wieloryby to rzadki gatunek, który od dawna budzi zainteresowanie naukowców i badaczy oceanów. Artykuł przedstawia ich unikalne cechy, w tym nietypowy kolor skóry, który nadaje im niepowtarzalny wygląd. Opisuje również ich imponujące rozmiary i specyfikę środowiska, w jakim żyją.
\end{abstract}

\part{"Tajemnicze Kruki Oceanu: Wprowadzenie do Świata Czarnych Wielorybów"}
\section{Czarne Wieloryby}\label{marker}
Od lat morza kryją w sobie niewyjaśnione tajemnice, a jednym z najbardziej enigmatycznych mieszkańców głębin jest czarny wieloryb. Jego nazwa nie tylko odnosi się do intensywnego koloru skóry, ale także do mrocznego aura, jaka otacza te majestatyczne stworzenia oceaniczne. W tym pierwszym rozdziale wyruszamy w fascynującą podróż ku poznaniu świata tych tajemniczych olbrzymów, zaczynając od ich unikalnych cech.
\subsection{Wizja Mroku Oceanicznego}
Zaledwie kropla w bezmiarze oceanów, budzą niekwestionowane zainteresowanie badaczy morskich. Głównym powodem jest ich niezwykły wygląd, który wyróżnia się spośród innych gatunków wielorybów. Skóra tych olbrzymów przybiera głęboki, niemal czarny odcień, nadając im niepowtarzalny charakter. Woda wokół nich staje się jak płótno, na którym malowana jest tajemnica nieznanego świata.
\newpage
\section{Niezwykła Kolorystyka Skóry}
Co sprawia, że skóra czarnych wielorybów jest tak nietypowa? To pytanie pozostaje otwarte, gdyż dotychczasowe badania nie dostarczyły jednoznacznej odpowiedzi. Spekulacje naukowe koncentrują się na możliwym związku między ich barwą a specyficznymi warunkami środowiska, w jakim żyją. Jednak, mimo że naukowcy zgłębiają tę zagadkę, kolorowa paleta czarnych wielorybów pozostaje jednym z najbardziej niezbadanych obszarów w oceanologii.
\subsection{W Poszukiwaniu Gigantów Oceanów}
Próba zrozumienia czarnych wielorybów to jednocześnie poszukiwanie olbrzymów oceanów. Te masywne stworzenia osiągają imponujące rozmiary, co czyni je jednymi z największych zwierząt na naszej planecie. Jednak, mimo ich potęgi, czarne wieloryby rzadko pojawiają się na powierzchni, co utrudnia badanie ich zachowań i migracji.
\section{W Cieniu Morskich Głębin}
Czarne wieloryby to istoty, które w większości czasu spędzają w ukryciu przed ludzkim okiem. Ich środowisko to głębiny oceanów, gdzie nieprzenikniony mrok staje się dla nich schronieniem. Badacze przy użyciu zaawansowanych technologii próbują odkryć sekrety ich naturalnego habitatu, a towarzyszące im tajemnicze dźwięki stają się przedmiotem badań akustyków morskich.
\subsection{Wędrówki przez Oceaniczne Krainy}
Podążając za migracjami czarnych wielorybów, naukowcy odkryli, że te potężne istoty przemierzają ogromne odległości, pozostawiając za sobą setki kilometrów oceanicznych przestrzeni. Co sprawia, że wyruszają w tak dalekie podróże? Czy to poszukiwanie pożywienia, czy też spełnienie innego biologicznego czy instynktownego celu? Odpowiedzi na te pytania są kluczowe dla zrozumienia roli, jaką odgrywają w ekosystemie morskim.
\newpage
\section{Słuchając Oceanicznych Pieśni}
Akustyka morska staje się narzędziem w rękach badaczy, umożliwiając im odsłuchanie pieśni oceanu. Czarne wieloryby komunikują się za pomocą charakterystycznych dźwięków, które mogą pełnić funkcje zarówno społeczne, jak i nawigacyjne. Odkrycie znaczenia tych dźwięków stawia pytanie, czy są one wyrazem inteligentnej komunikacji czy też jedynie instynktownych reakcji na otaczające środowisko.
\subsection{Akustyka morska}
Staje się narzędziem w rękach badaczy, umożliwiając im odsłuchanie pieśni oceanu. Czarne wieloryby komunikują się za pomocą charakterystycznych dźwięków, które mogą pełnić funkcje zarówno społeczne, jak i nawigacyjne. Odkrycie znaczenia tych dźwięków stawia pytanie, czy są one wyrazem inteligentnej komunikacji czy też jedynie instynktownych reakcji na otaczające środowisko.
\newline
\newline
\newline
\textbf{W kolejnych rozdziałach artykułu wnikniemy głębiej w tajemnice życia czarnych wielorybów, przyjrzymy się ich ekologii, zachowaniom społecznym i zagrożeniom, jakie mogą być dla nich wynikiem zmian w środowisku morskim. Celem tej podróży jest odkrycie, co sprawia, że czarne wieloryby są nie tylko jednymi z najbardziej uderzających widoków w oceanach, ale także kluczowymi graczami w delikatnej równowadze ekosystemu morskiego.}
\newpage
\part{"Tajemnice Czarnej Pustki: Gdzie Czarni Wieloryby Przemierzają Oceaniczne Otchłanie"}
\textsl{W drugim rozdziale przenosimy się w głąb oceanicznych otchłani, gdzie czarne wieloryby odnajdują swoje naturalne środowisko. Rozważymy ich migracje, tajemnicze szlaki, oraz to, co sprawia, że te potężne stworzenia decydują się na wędrówki po głębinach mórz.}
\section{Życie Rodzinne Czarnych Wielorybów}
Czarne wieloryby, choć w pierwszym rzucie oka mogą jawić się jako samotnicy, zaskakują swoimi głębokimi więziami rodzinno-społecznymi. W rodzinach tych olbrzymów panuje hierarchia, a relacje między członkami są złożone. Matki odgrywają kluczową rolę w opiece nad potomstwem, a młode wieloryby uczą się od starszych członków grupy.
\subsection{Taniec Wielorybich Duetów}
Jednym z najbardziej fascynujących aspektów społecznych czarnych wielorybów jest ich zdolność do wykonywania skomplikowanych i synchronicznych tańców. Naukowcy zastanawiają się, czy te widowiskowe widowiska mają funkcje komunikacyjne, czy też pełnią inne, bardziej głębokie znaczenia w kontekście społeczności wielorybowej.
\newpage
\section{Solidarność Wielorybich Społeczności}
W trudnych warunkach oceanicznych solidarność społeczności czarnych wielorybów staje się kluczowym elementem przetrwania. Wspólne polowania, obrona przed drapieżnikami czy nawet pomoc w przypadku choroby jednego z członków grupy, to przykłady zachowań, które podkreślają, że te wieloryby żyją w społeczeństwach opartych na współpracy i wzajemnej pomocy.
\paragraph{Mowa Wielorybów jako Zapis Matematyczny}
W miarę jak badania nad komunikacją wielorybów stają się coraz bardziej zaawansowane, naukowcy odkrywają, że "mowa" tych morskich olbrzymów może być porównywana do precyzyjnego zapisu matematycznego. W tej sekcji przyjrzymy się fascynującym aspektom języka wielorybów, jego strukturze i ewolucji, które sugerują, że istnieje coś więcej niż jedynie serenada oceanicznych dźwięków.
\newline
\large{ODCZYTANE ZDANIA WIELORYBOW : }\\
\\\\
$x^{n + 2} = \lambda + \int_{1}^{\infty} \frac{\sum_{i=1}^{n} y_{i}}{\sqrt[3]{2 + z_{j}}} dz.$
\\\\
$ \frac{y_{i}}{\sqrt[3]{2 + z_{j}}}.$ \footnotesize{      TO PODOBNO ZNACZY}"TAK"
\\\\\\\\
$\biggl(\sum_{n\ge 1} a_{n} n^{-s}\biggr)^2$


\begin{figure}[h]
  \centering
  \begin{tikzpicture}
    \draw (2,1) circle (0.8);
  \end{tikzpicture}
  \caption{kolko \cite{jeden}}
  \label{fig:rysunek}
\end{figure}

\begin{figure}[h]
  \centering
  \begin{tikzpicture}
    \draw (0,0) rectangle (4,2);
    \draw (2,1) node[text width=3cm,align=center] {CZARNY WIELORYB};
  \end{tikzpicture}
  \caption{kwadrat z tekstem \cite{dwa}}
  \label{fig:rysunek-tekst}
\end{figure}

\newpage
\section{Tabela}

Poniżej znajduje się tabela, która odnosi się do tekstu w artykule z sekcji \ref{marker} cz. 1.

\begin{table}[h]
  \centering
  \begin{tabular}{|c|c|}
    \hline
    \textbf{Rodzaj} & \textbf{Opis} \\
    \hline
    Czarny Wieloryb & szybki silny fajny \\
    Wieloryb & nudny zwykly \\
    Pletwal Blekitny & najwiekszy wolny \\
    \hline
  \end{tabular}
  \caption{Odniesienie do \ref{marker}. Czarne Wieloryby}
  \label{tab:odniesienie}
\end{table}

\section{Podsumowanie}

To jest podsumowanie mojego artykułu. \\
Odniesienie do tabeli : \ref{tab:odniesienie}.\\
Odniesienie do rysunku bez tekstu : Rysunek \ref{fig:rysunek}\\
Odniesienie do rysunku z tekstem : Rysunek \ref{fig:rysunek-tekst}\\
\bibliographystyle{plain}
\bibliography{mamdosc.bib}


\end{document}